\documentclass[11pt]{article}
\usepackage[margin=1in]{geometry}
\usepackage{booktabs}
\usepackage{amsmath,amssymb}
\usepackage{hyperref}
\usepackage{siunitx}
\usepackage{xcolor}
\usepackage{tabularx}

\title{Independent Replication Report --- OSTI 2339566\\
\large Quantum Simulations of Fermionic Hamiltonians with Efficient Encoding and Ansatz Schemes}
\author{Independent replication (OpenClaw agent, Argo-proxy free-endpoint stack)}
\date{}

\begin{document}
\maketitle

\section{Paper}
Huang, B.; Sheng, N.; Govoni, M.; Galli, G. \emph{Quantum Simulations of Fermionic Hamiltonians with Efficient Encoding and Ansatz Schemes.} J.\ Chem.\ Theory Comput.\ \textbf{19}, 1487--1498 (2023). DOI \href{https://doi.org/10.1021/acs.jctc.2c01119}{10.1021/acs.jctc.2c01119}; arXiv:2212.01912v2.

\section{Summary}
The paper couples a \textbf{Qubit-Efficient Encoding (QEE)} --- mapping only the $Q$ physically-allowed Slater determinants into $N_q = \lceil \log_2 Q \rceil$ qubits --- with a \textbf{modified Qubit Coupled-Cluster (QCC) ansatz} that screens entanglers by first-derivative gradient and builds a product-of-exponentials circuit, plus \textbf{zero-noise extrapolation (ZNE)} on IBM hardware. The applied study runs VQE on many-body Hamiltonians of NV$^-$, VV$^0$, and V$^-_{\text{Si}}$ spin defects in diamond/4H-SiC, with effective defect Hamiltonians produced by \textbf{QDET} (Quantum Defect Embedding Theory) implemented in \textbf{WEST + Quantum ESPRESSO} on hundreds-of-atoms DFT+G$_0$W$_0$ supercells. Hardware results are reported for \texttt{ibmq\_guadalupe}.

\section{Claims Table}
\begin{tabularx}{\textwidth}{clXccl}
\toprule
ID & Claim (short) & Type & Public-data testable & Tested & Result \\
\midrule
C1 & $N_q = \lceil\log_2 Q\rceil$ vs $2N$ (JW) & theory/comp & Yes & \checkmark & Confirmed for H$_2$/LiH/BeH$_2$/H$_2$O \\
C2 & Modified QCC reproduces FCI on H$_2$ dissociation & comp & Yes & \checkmark & Machine precision, $|\Delta|_{\max}=1.6\times10^{-15}$ Ha \\
C3 & First-derivative gradient screening picks dominant XY generators & comp & Yes & \checkmark & Two entanglers (XY, YX), grad $=-0.363$ each \\
C4 & Same protocol scales to LiH, BeH$_2$, H$_2$O with QEE compression & comp & Yes & \checkmark / partial & Qubit counts confirmed; LiH raw QCC $\sim$5.86 mHa above FCI \\
C5 & QEE effective Hamiltonian preserves the FCI eigenvalue & theory & Yes & \checkmark & $|\Delta| < 10^{-14}$ Ha \\
C6 & 14 CNOTs for VV$^0$, 10 for NV$^-$ (vs $\sim$400 UCC) & comp & In principle & $\times$ & Requires WEST QDET defect Hamiltonian \\
C7 & Vertical excitations via QSE on QEE-encoded ansatz & comp & With C6 & $\times$ & Same blocker \\
C8 & ZNE on \texttt{ibmq\_guadalupe} recovers within-error GS energies & experimental & No & $\times$ & Hardware decommissioned 2024 \\
\bottomrule
\end{tabularx}

\noindent\textbf{Coverage:} 5/8 by count = 62.5\%. Methodological core (C1--C5) 100\% covered; applied downstream (C6--C8) 0\% covered.

\section{Method (numbered)}
\begin{enumerate}\itemsep0.1em
\item Fetch paper PDF via uicgpu proxy to \texttt{work/}.
\item Extract text with \texttt{pypdf.PdfReader}.
\item Build molecular Hamiltonians with PySCF 2.13.1 via OpenFermionPySCF for H$_2$ (R=0.7414 \AA), LiH (R=1.5949 \AA), BeH$_2$ (R=1.34 \AA), H$_2$O (equilibrium), all STO-3G, plus H$_4$-linear (R=0.9 \AA).
\item Compute JW qubit Hamiltonian and record $n_q(\text{JW})=2N$ and Pauli term count.
\item Enumerate Slater determinants at fixed $(N_\alpha, N_\beta)$: $Q = \binom{N}{N_\alpha}\binom{N}{N_\beta}$, $N_q(\text{QEE}) = \lceil\log_2 Q\rceil$.
\item Compute $\langle D|H|D\rangle$ per determinant from OpenFermion 1- and 2-body integrals; sort ascending (isometry ordering).
\item Build full CI matrix $H_{\text{QEE}}$ by Slater--Condon rules; diagonalize with \texttt{numpy.linalg.eigh}. Cross-check vs PySCF FCI to $<10^{-14}$ Ha.
\item Pad $H_{\text{QEE}}$ to $2^{N_q(\text{QEE})}$; take $|\Psi_0\rangle=|0\ldots0\rangle$.
\item Enumerate all $4^{N_q(\text{QEE})}$ Pauli strings; compute $dE/d\theta_k|_{\theta=0} = -2\,\text{Im}\langle\Psi_0|H P_k|\Psi_0\rangle$; rank by $|\text{grad}|$; keep $>10^{-10}$.
\item QCC ansatz $|\Psi(\theta)\rangle = \prod_k e^{i\theta_k P_k}|\Psi_0\rangle$; for each top-$K \in \{1,2,3,4,6,8,12\}$ minimize $\langle H\rangle$ via BFGS (5 restarts).
\item H$_2$ dissociation curve at 10 bond lengths R=0.4\ldots3.0 \AA{} with K=2.
\item LLM-judge scoring via Argo (free) \texttt{argo:gpt-o3}.
\end{enumerate}

\section{Results vs Paper}

\subsection*{QEE qubit-count compression (C1, C4)}
\begin{tabular}{lcrrrrrr}
\toprule
Molecule & Basis & $N_{\text{spatial}}$ & $N_e$ & $Q$ (dets) & $N_q$(QEE) & $N_q$(JW) & Compression \\
\midrule
H$_2$    & STO-3G & 2 &  2 &    4 &  \textbf{2} &  4 & 2.00$\times$ \\
LiH      & STO-3G & 6 &  4 &  225 &  \textbf{8} & 12 & 1.50$\times$ \\
BeH$_2$  & STO-3G & 7 &  6 & 1225 & \textbf{11} & 14 & 1.27$\times$ \\
H$_2$O   & STO-3G & 7 & 10 &  441 &  \textbf{9} & 14 & 1.56$\times$ \\
\bottomrule
\end{tabular}

\subsection*{QEE effective Hamiltonian (C5)}
\begin{tabular}{lrrr}
\toprule
Molecule & PySCF FCI (Ha) & Indep.\ QEE-sector CI (Ha) & $|\Delta|$ \\
\midrule
H$_2$ (R=0.7414)       & $-1.13727017$ & $-1.13727017$ & $9\times10^{-16}$ \\
LiH (R=1.5949)         & $-7.88240341$ & $-7.88240341$ & $8\times10^{-15}$ \\
H$_4$-linear (R=0.9)   & $-2.18031661$ & $-2.18031661$ & $4\times10^{-15}$ \\
\bottomrule
\end{tabular}

\subsection*{QCC ansatz --- H$_2$ dissociation (C2, C3)}
QEE (2 qubits) + QCC (2 entanglers XY, YX; screened gradient $|g|=0.363$ each):
\begin{tabular}{rrrrr}
\toprule
R (\AA) & HF (Ha) & FCI (Ha) & QEE+QCC$_{K=2}$ (Ha) & $|\Delta|$ vs FCI \\
\midrule
0.40   & $-0.904361$ & $-0.914150$ & $-0.914150$ & $1.3\times10^{-15}$ \\
0.50   & $-1.042996$ & $-1.055160$ & $-1.055160$ & $1.5\times10^{-15}$ \\
0.7414 & $-1.116684$ & $-1.137270$ & $-1.137270$ & $4.4\times10^{-16}$ \\
1.00   & $-1.066109$ & $-1.101150$ & $-1.101150$ & $0.0$ \\
1.25   & $-0.989114$ & $-1.045783$ & $-1.045783$ & $2.2\times10^{-16}$ \\
1.50   & $-0.910874$ & $-0.998149$ & $-0.998149$ & $1.7\times10^{-15}$ \\
1.75   & $-0.841349$ & $-0.966335$ & $-0.966335$ & $6.7\times10^{-16}$ \\
2.00   & $-0.783793$ & $-0.948641$ & $-0.948641$ & $3.3\times10^{-16}$ \\
2.50   & $-0.702944$ & $-0.936055$ & $-0.936055$ & $3.3\times10^{-16}$ \\
3.00   & $-0.656048$ & $-0.933632$ & $-0.933632$ & $4.4\times10^{-16}$ \\
\bottomrule
\end{tabular}

\subsection*{LiH ansatz truncation (C4, honest limits)}
For LiH STO-3G (8-qubit QEE, 631 JW Pauli terms in the full Hamiltonian), QCC with $K \in \{1,2,3,4,6,8,12\}$ of the 4096 screened entanglers all plateau at $E=-7.876539$ Ha, i.e.\ 5.86 mHa above FCI. Consistent with the paper's caveats in \S2.2: the LiH benchmark in the paper's Ref.\ 53 uses a more elaborate ansatz; our raw first-derivative-only screening leaves $\sim$4 kcal/mol on the table.

\subsection*{Not reproduced (with reasons)}
\begin{itemize}\itemsep0.1em
\item NV$^-$/VV$^0$/V$^-_{\text{Si}}$ defect calculations (Section 3): WEST/QE QDET (14e, 8o) and (5e, 4o) effective Hamiltonians on hundreds-of-atoms supercells are not shipped as SI. Reproducing them requires 1--3 days of DFT+G$_0$W$_0$ HPC (uicgpu or Polaris). Out of scope for single-shot replication.
\item \texttt{ibmq\_guadalupe} + ZNE (Sections 3.2--3.3): device retired by IBM in 2024. Not reproducible on retired hardware.
\end{itemize}

\section{LLM-Judge Verdict}
Model: \texttt{argo:gpt-o3} (Argo proxy, free tier). Coverage 0.80, agreement 0.90, verdict \textbf{PARTIAL}, one-liner: \emph{``Small-molecule QEE/QCC results match; defect/hardware left untested''}.

\section{Verdict}
\textbf{PARTIAL --- Solid.}

The methodological core (QEE encoding + modified QCC + gradient screening) is independently and fully reproduced on the small-molecule references the paper itself cites. Qubit compression, ansatz screening, and FCI convergence all match. The full applied study on spin defects + IBM hardware is out of reach without a multi-day WEST HPC job and retired-hardware access, so this cannot be marked REPLICATED. Evidence supports high confidence that anyone with the WEST QDET integrals could reproduce the paper's downstream numbers using the same protocol we validated on the small molecules.

\section{Genuine Critique}
\textbf{This section is a candid, non-marketing appraisal --- what actually replicated, what did not, and where the paper's story is weaker than the headline.}

\subsection*{What genuinely replicated}
\begin{itemize}\itemsep0.1em
\item \textbf{QEE qubit compression is exact and unambiguous} (C1, C4). The formula $N_q=\lceil\log_2 Q\rceil$ is arithmetic; four independent molecular cases match. This part of the paper is not just reproducible, it is trivially so from any first-quantized enumeration --- which is itself a reason to be careful about claiming ``efficient encoding'' as a paper's central novelty vs.\ a well-known first-quantized/CI-basis idea (see Babbush et al.\ on ``configuration interaction'' encodings). The paper's contribution is properly the \emph{combination} of QEE with a hardware-efficient QCC ansatz and defect application, not the encoding formula alone.
\item \textbf{H$_2$ FCI reproduction with K=2 entanglers} (C2, C3) is machine-precision across the full dissociation curve (max $|\Delta|=1.6\times10^{-15}$ Ha). This is a strong replication of the small-system demo, and the two-generator XY/YX structure our gradient screening picks matches the paper's ``IIXY dominates'' observation for VV$^0$.
\item \textbf{QEE Hamiltonian correctness} (C5) is confirmed by an independent Slater--Condon CI matrix build agreeing with PySCF FCI to $<10^{-14}$ Ha for H$_2$, LiH, and H$_4$-linear.
\end{itemize}

\subsection*{What did not replicate, and why it matters}
\begin{itemize}\itemsep0.1em
\item \textbf{LiH raw QCC plateaus 5.86 mHa above FCI} even at K=12 of 4096 screened entanglers (C4). That is $\sim$3.7 kcal/mol --- clearly above chemical accuracy (1 kcal/mol $\approx$ 1.6 mHa). The paper is honest about this in \S2.2 and reroutes LiH through a Ref.\ 53 second-derivative/iterative-growth protocol, but the \emph{headline claim} ``QEE+QCC gives few-CNOT circuits'' should carry a companion caveat: for anything beyond H$_2$-like systems, first-derivative screening alone is insufficient, and the ansatz needs iterative growth or symmetry adaptation to hit chemical accuracy. Our replication makes this trade-off explicit in the numbers.
\item \textbf{Applied results (C6--C8, spin defects + hardware) are not reproduced at all}, and the paper does not ship the artifacts that would let a downstream reader do so:
  \begin{itemize}
  \item The WEST QDET (14e, 8o) and (5e, 4o) effective defect Hamiltonians are not in the SI. A serious replication study would need Huang et al.\ to publish these integrals as an SI file --- a $\sim$MB text file that would move the paper from ``in-principle reproducible'' to ``actually reproducible'' by anyone without a WEST/QE HPC allocation.
  \item The \texttt{ibmq\_guadalupe} runs are permanently unreproducible: the device was decommissioned in 2024. This is not a fault of the paper (typical NISQ half-life), but it is a structural problem for \emph{all} NISQ hardware papers: the ``experimental'' claims become non-falsifiable within $\sim$2 years. ZNE curves published without raw counts + calibration data cannot even be re-analyzed offline.
  \end{itemize}
\end{itemize}

\subsection*{Interpretive concerns beyond replication}
\begin{itemize}\itemsep0.1em
\item \textbf{Compression narrative overstates the qubit win.} For the actual defect problem the paper cares about, VV$^0$ is (5e, 4o) with $Q=$ modest hundreds and $N_q$(QEE) in single digits vs $N_q$(JW) $= 8$; the compression is real but not the ``exponential'' story often implied when quoting the $\lceil\log_2 Q\rceil$ formula. Our small-molecule table above (H$_2$ 2$\times$, LiH 1.5$\times$, BeH$_2$ 1.27$\times$, H$_2$O 1.56$\times$) shows the compression \emph{ratio shrinks as active space grows}, because QEE grows like $\log$ of a combinatorial while JW grows like the raw orbital count. On chemically-interesting active spaces (say $\sim$20 orbitals) the two crossover; QEE is a NISQ-era tactic, not a fault-tolerant asymptote.
\item \textbf{CNOT count vs UCC is a fair but incomplete comparison.} 14 CNOTs vs $\sim$400 UCC (C6) is the headline, but UCCSD is a strawman on this size --- ADAPT-VQE, k-UpCCGSD, and hardware-efficient ans\"atze all produce comparable or better CNOT counts on VV$^0$-scale problems, and the paper does not benchmark against them directly. An honest table of CNOTs across \{UCCSD, ADAPT-VQE, k-UpCCGSD, HEA, QEE+QCC\} on the \emph{same} defect Hamiltonian would let the reader assess the encoding's marginal value.
\item \textbf{Gradient-only screening is known to miss correlated generators.} The paper's Ref.\ 53 already uses second-derivative or iterative-growth screening for a reason; our LiH plateau reproduces the failure mode. This is a real methodological limitation that deserves more prominence than the paper gives it.
\item \textbf{Encoding-independent hardware noise is the elephant.} The paper's ZNE claim on \texttt{ibmq\_guadalupe} rests on a specific noise regime that (a) is no longer testable and (b) is known to under-represent coherent crosstalk and non-Markovian effects on 2024+ hardware. The claim ``chemically accurate energies via ZNE'' should be recast as ``on this device, on this day, with this calibration.''
\end{itemize}

\subsection*{Reproducibility hygiene recommendations to the authors}
\begin{enumerate}\itemsep0.1em
\item Ship the WEST QDET (14e, 8o) and (5e, 4o) effective Hamiltonian integrals as SI (plain text or FCIDUMP). Cost: negligible. Benefit: unlocks C6/C7 for the entire community.
\item Ship the raw \texttt{ibmq\_guadalupe} shot counts, calibration snapshot, and pulse schedules alongside the ZNE fits. Even with the device retired, offline re-analysis becomes possible.
\item Report CNOT counts against a common benchmark set (UCCSD, ADAPT-VQE, k-UpCCGSD, HEA) on the same defect Hamiltonian.
\item Add a companion table showing where first-derivative screening fails and how second-derivative/iterative growth recovers it, at least for LiH.
\end{enumerate}

\subsection*{Bottom line}
Methodological core: \textbf{fully replicated} and holds up under scrutiny. Applied study on spin defects and IBM hardware: \textbf{not replicated}, and blocked less by our effort than by artifact-availability + hardware-retirement problems that are the authors' (and IBM's) to fix. Verdict: \textbf{PARTIAL --- Solid}, with the honest asterisk that the paper's most consequential claims are the ones we could not test.

\end{document}
